% =============================================================================
% ForensiQ — Proposta de Projecto
% UC 21184 · Projeto de Engenharia Informática · UAb
% =============================================================================
\documentclass[a4paper,11pt]{article}

\usepackage[utf8]{inputenc}
\usepackage[T1]{fontenc}
\usepackage[portuguese]{babel}
\usepackage[top=2.5cm, bottom=2.5cm, left=3cm, right=2.5cm]{geometry}
\usepackage{booktabs}
\usepackage{longtable}
\usepackage{array}
\usepackage{parskip}
\usepackage{hyperref}
\usepackage{enumitem}

\hypersetup{
    colorlinks=true,
    linkcolor=black,
    urlcolor=blue,
    pdfauthor={João Rodrigues},
    pdftitle={ForensiQ — Proposta de Projecto}
}

% Colunas de largura fixa para tabelas
\newcolumntype{L}[1]{>{\raggedright\arraybackslash}p{#1}}

% =============================================================================
\begin{document}
% =============================================================================

\begin{center}
    {\LARGE\bfseries Proposta de Projecto}\\[0.8em]
    {\large ForensiQ --- Plataforma Modular de Gestão de Prova Digital\\
    para First Responders}\\[1.2em]
    \begin{tabular}{ll}
        \textbf{Estudante:}   & João Rodrigues · 2203474 \\
        \textbf{Orientador:}  & Pedro Pestana \\
        \textbf{Data:}        & 22 de março de 2026 \\
        \textbf{Versão:}      & 1.0 \\
    \end{tabular}
\end{center}

\vspace{0.5em}
\hrule
\vspace{1.5em}

% =============================================================================
\section{Sinopse}
% =============================================================================

Os agentes que chegam primeiro a uma cena de crime --- os \emph{first responders} ---
recorrem actualmente a métodos manuais (papel, fotografias avulsas, anotações) para
registar e preservar a prova digital. Esta ausência de padronização traduz-se em
vulnerabilidades processuais reais: cadeia de custódia inconsistente, prova sem
localização verificável e procedimentos que variam de agente para agente. O problema
é identificado da experiência directa do estudante como perito digital forense na PSP.

O ForensiQ é uma aplicação web \emph{mobile-first} que digitaliza e padroniza este
processo, do momento da apreensão no terreno até à chegada ao laboratório e remessa
das conclusões da perícia. Distingue-se dos sistemas existentes (SAFE, Axon Evidence)
por ser desenhado para o agente sem formação técnica avançada, por seguir os
procedimentos da PSP, e por garantir integridade dos registos via hash SHA-256
conforme a ISO/IEC 27037. Diferencia-se do Projecto \#38 (OSINT) por não adquirir
prova --- apenas a gerir e controlar após apreensão.

O resultado esperado é um protótipo funcional demonstrável num smartphone no terreno,
com API REST documentada, módulo de forense digital completo e cadeia de custódia
com log imutável verificável. O sucesso é medido pelos critérios de aceitação
definidos para cada funcionalidade do MVP.

% =============================================================================
\section{MVP --- Definição e critérios de aceitação}
% =============================================================================

\subsection{Funcionalidade 1 --- Autenticação com perfis (agente e perito)}

\textbf{Critério de aceitação:}
Dado email e password válidos de agente, o sistema autentica e devolve token JWT em
menos de 2 segundos; dado email inválido, devolve erro 401 sem expor informação
interna. Perfis agente e perito têm permissões distintas verificáveis.

\subsection{Funcionalidade 2 --- Registo de prova com fotografia e metadados GPS}

\textbf{Critério de aceitação:}
Um agente autenticado cria um registo de prova com fotografia, tipo, descrição e
localização GPS capturada automaticamente pelo browser. O registo recebe
identificador único, timestamp do servidor e hash SHA-256 dos metadados.
Verificável em menos de 30 segundos num telemóvel.

\subsection{Funcionalidade 3 --- Cadeia de custódia com máquina de estados e log imutável}

\textbf{Critério de aceitação:}
A prova percorre estados em ordem fixa sem possibilidade de retrocesso:
\textsc{apreendida} $\rightarrow$ \textsc{em transporte} $\rightarrow$
\textsc{recebida no laboratório} $\rightarrow$ \textsc{em perícia} $\rightarrow$
\textsc{concluída} $\rightarrow$ \textsc{devolvida/destruída}.
Cada transição regista agente responsável, timestamp e hash do registo nesse
momento. Nenhuma transição pode ser saltada. O log é \emph{append-only} ---
sem \texttt{UPDATE} nem \texttt{DELETE} na tabela de custódia. Qualquer
tentativa de alteração directa na base de dados é detectável via hash.

\subsection{Funcionalidade 4 --- Módulo de forense digital (ficha de dispositivo)}

\textbf{Critério de aceitação:}
Perito preenche ficha completa de dispositivo digital (tipo, marca, modelo, estado,
número de série) associada a uma prova existente.

\subsection{Funcionalidade 5 --- API REST documentada}

\textbf{Critério de aceitação:}
Mínimo 10 endpoints funcionais documentados e testáveis via Swagger UI
(\texttt{drf-spectacular}).

\subsection{Funcionalidade 6 --- Exportação de relatório de ocorrência em PDF}

\textbf{Critério de aceitação:}
Relatório gerado automaticamente com todos os registos de prova e histórico de
custódia da ocorrência.

% =============================================================================
\section{Stack tecnológica}
% =============================================================================

\begin{longtable}{L{3cm} L{5cm} L{6.5cm}}
    \toprule
    \textbf{Componente} & \textbf{Tecnologia escolhida} & \textbf{Justificação} \\
    \midrule
    \endhead
    Backend
        & Django + Django REST Framework
        & Estrutura convencional e previsível; autenticação built-in; ORM que escreve SQL \\[0.4em]
    Frontend
        & HTML + CSS + JavaScript vanilla
        & Mobile-first sem overhead de framework; suficiente para o MVP \\[0.4em]
    Base de dados
        & PostgreSQL
        & Integridade referencial crítica para cadeia de custódia; append-only para logs \\[0.4em]
    Autenticação
        & JWT + RBAC
        & Dois perfis: agente (first responder) e perito forense digital \\[0.4em]
    API docs
        & drf-spectacular $\rightarrow$ Swagger UI
        & Documentação automática da API REST \\[0.4em]
    Geolocalização
        & Browser Geolocation API + Leaflet.js
        & Sem dependência de API externa paga; funciona no telemóvel \\[0.4em]
    Integridade
        & SHA-256 (\texttt{hashlib})
        & Hash dos metadados do registo no momento de criação \\[0.4em]
    CI/CD
        & GitHub Actions
        & Pipeline de testes automáticos a cada push \\
    \bottomrule
\end{longtable}

% =============================================================================
\section{Esboço de arquitectura --- C4 Nível 1}
% =============================================================================

\textbf{Sistema:} ForensiQ

\textbf{Utilizadores:}
\begin{itemize}[nosep]
    \item Agente (first responder) --- regista provas no terreno, inicia cadeia de
    custódia, captura fotografias e localização GPS.
    \item Perito forense digital --- recebe provas no laboratório, preenche fichas
    de dispositivo, avança estados da custódia, exporta relatórios.
\end{itemize}

\textbf{Sistemas externos:}
\begin{itemize}[nosep]
    \item Browser Geolocation API --- fornece coordenadas GPS do dispositivo do
    agente no terreno.
    \item Servidor de email (futuro) --- notificações de transição de custódia.
\end{itemize}

% =============================================================================
\section{Calendário individual detalhado}
% =============================================================================

\begin{longtable}{L{1.8cm} L{3.2cm} L{7cm} L{3cm}}
    \toprule
    \textbf{Semanas} & \textbf{Datas} & \textbf{Conteúdo planeado} & \textbf{Marco} \\
    \midrule
    \endhead
    Sem. 1--2  & 17--28 mar
        & Proposta completa. Configuração do repositório. Sinopse, MVP e calendário ao orientador.
        & \textbf{Proposta (25 mar)} \\[0.4em]
    Sem. 3--4  & 31 mar--11 abr
        & MoSCoW final, diagramas C4, modelo de dados ER, ADRs obrigatórios.
        & \\[0.4em]
    Sem. 5--6  & 14--25 abr
        & Wireframes mobile-first. Início implementação: modelos Django, autenticação JWT, registo de prova.
        & \\[0.4em]
    Sem. 7     & 28 abr--2 mai
        & Demo interna ao orientador (Teams). Cadeia de custódia funcional.
        & \textbf{Demo interna} \\[0.4em]
    Sem. 8     & 5--6 mai
        & Relatório intercalar.
        & \textbf{Intercalar (6 mai)} \\[0.4em]
    Sem. 9--10 & 7--16 mai
        & Módulo forense digital. API REST completa (10+ endpoints). Swagger UI.
        & \\[0.4em]
    Sem. 11--12 & 19--30 mai
        & Exportação PDF. Testes integração (Postman/Newman). GitHub Actions CI.
        & \\[0.4em]
    Sem. 13    & 2--6 jun
        & Polimento UI. Revisão relatório final (Cap. 4 e 5).
        & \\[0.4em]
    Sem. 14    & 9--13 jun
        & Revisão final do relatório. Correcções de última hora.
        & \\[0.4em]
    Sem. 15    & 16--20 jun
        & Preparação da defesa oral (15 min). Ensaio com simulação de perguntas.
        & \textbf{Prep. defesa} \\[0.4em]
    Sem. 16    & 24 jun
        & Submissão do relatório final.
        & \textbf{Final (24 jun)} \\
    \bottomrule
\end{longtable}

\end{document}
